\chapter{Conclusiones}
\label{chap:conclusiones}

\section{Respuesta a la hipótesis}

La hipótesis se sostiene dentro del alcance definido: una simulación de eventos discretos permite convertir el diseño de un banco QKD en preguntas medibles, cuantificar la variación entre corridas y detectar cuándo una salida exige revisar sus supuestos. El trabajo no confirma la factibilidad física del Campus Miguelete ni certifica seguridad criptográfica.

Los cuatro criterios de verificación se cumplen. La evidencia conserva parámetros, semillas, bits, errores, tiempos y clics; cada punto contiene \PThreeRepetitions{} repeticiones; el flujo comprueba invariantes numéricos; y la referencia analítica detecta una pérdida de holgura que la lectura aislada de QBER no habría mostrado. Este último resultado es importante: la utilidad de la simulación no consiste sólo en producir curvas plausibles, sino también en impedir que una tasa incoherente se transforme en especificación de diseño.

\section{Cumplimiento de objetivos}

\begin{table}[htbp]
  \centering
  \small
  \begin{tabularx}{\textwidth}{p{0.27\textwidth}p{0.18\textwidth}X}
    \toprule
    Objetivo & Estado & Evidencia \\
    \midrule
    Escenario Alice--Bob \textit{time-bin} & Cumplido & SeQUeNCe 0.8.5, componentes, canales y parche identificados \\
    Barrido de distancia & Cumplido con límite & Transición de control entre \num[round-mode=places,round-precision=2]{\PThreeTransitionLowKm} y \SI[round-mode=places,round-precision=2]{\PThreeTransitionHighKm}{\kilo\meter}; no es alcance físico \\
    Sensibilidad del detector & Cumplido, no concluyente & Los IC de eficiencia y fondo incluyen cero \\
    Sensibilidad a visibilidad & Cumplido internamente & Pendiente QBER negativa, con doce puntos sin holgura de tasa \\
    Comparación señuelo & Cumplido como estimador & Igual \(\mu\), cota asintótica híbrida y reemplazos conservadores registrados \\
    Diseño Campus Miguelete & Parcial & Topología y criterios definidos; faltan ruta y pérdidas medidas \\
    Trazabilidad & Cumplido & 60 agregados, 1.800 registros, 2.100 ejecuciones, macros y hashes \\
    \bottomrule
  \end{tabularx}
  \caption{Respuesta de la evidencia a los objetivos del trabajo.}
  \label{tab:objetivos_cierre}
\end{table}

\section{Conclusiones por experimento}

El barrido de distancia reproduce la caída de ganancia, pero su conclusión principal no es una distancia máxima. Entre \num[round-mode=places,round-precision=2]{\PThreeTransitionLowKm} y \SI[round-mode=places,round-precision=2]{\PThreeTransitionHighKm}{\kilo\meter}, la tasa tamizada deja de conservar holgura frente a la referencia. Los puntos posteriores sirven para auditar temporización y detecciones, no para defender prestaciones.

El experimento de detector no resuelve los efectos entre extremos. Para eficiencia, el intervalo de \(\Delta I\) va de \num[round-mode=places,round-precision=1]{\PThreeEffDifferenceLowKbps} a \num[round-mode=places,round-precision=1]{\PThreeEffDifferenceHighKbps}\,kbit/s y \(p=\num[round-mode=places,round-precision=3]{\PThreeEffPValue}\). Para fondo, va de \num[round-mode=places,round-precision=1]{\PThreeDarkDifferenceLowKbps} a \num[round-mode=places,round-precision=1]{\PThreeDarkDifferenceHighKbps}\,kbit/s y \(p=\num[round-mode=places,round-precision=3]{\PThreeDarkPValue}\). Esto impide ordenar detectores a partir de estas curvas. Una comparación de hardware debe incorporar ventanas, tiempo muerto, fluctuación temporal y pérdidas internas medidas.

El barrido de visibilidad confirma una pendiente negativa de QBER. Era esperable porque \(p_\mathrm{fase}\) se calcula desde \(V\); el resultado demuestra que la dependencia llega a BB84 después de corregir el interferómetro. No demuestra estabilidad física. Los doce puntos carecen de holgura de tasa, de modo que el valor absoluto del indicador no se usa como evidencia de factibilidad.

La comparación de estados señuelo muestra la utilidad de estimar la contribución monofotónica. Con igual \(\mu=0{,}1\), la tasa sin señuelos se anula cerca de \SI{24}{\kilo\meter}, mientras la estimación señuelo conserva \SI[round-mode=places,round-precision=1]{\PThreeDecoyTwentyFourKbps}{\kilo bit\per\second}. El resultado depende de ganancias analíticas, error de fase aproximado y tratamiento asintótico. Los \PThreeDecoyFallbackTotal{} reemplazos conservadores en 300 repeticiones pareadas deben considerarse al interpretar la media.

\section{Aporte para la implementación UNSAM}

El trabajo entrega una secuencia concreta para la siguiente etapa. Primero se mide una ruta óptica y se registra cada pérdida fija. Después se selecciona una combinación coherente de fuente, moduladores, detector, registro temporal e interferómetro. Esos valores restringen el dominio del simulador. Finalmente, se compara una medición corta con la salida del modelo y se ajustan temporización, fondo y visibilidad.

El presupuesto y las imágenes de equipos cumplen una función de preselección, no de compra. Los rangos de costo son orientativos y no provienen de cotizaciones vigentes. Antes de adquirir equipos se necesita una tabla con cantidad, modelo exacto, accesorios, precio, fecha, importación y contingencia. Sin ese paso, sumar mínimos y máximos produce un número aparente pero no un presupuesto ejecutable.

\section{Limitaciones}

Las principales limitaciones son:

\begin{itemize}
  \item la referencia de tasa no reproduce toda la temporización interna y detecta una discrepancia todavía no resuelta;
  \item la ventana de \SI{1}{\nano\second} es analítica y no una compuerta implementada en el receptor;
  \item el barrido de eficiencia excede la especificación del detector comercial usado como referencia;
  \item el modelo no separa errores por base ni ejecuta autenticación, reconciliación y amplificación de privacidad;
  \item el estimador señuelo es asintótico, híbrido y no componible;
  \item no se modelan ataques laterales, deriva térmica, dispersión, pérdidas internas ni calibración dinámica;
  \item el campus no cuenta aún con una ruta óptica medida ni una lista de materiales cotizada.
\end{itemize}

\section{Trabajo futuro priorizado}

La primera prioridad es resolver la diferencia entre tasa tamizada y referencia. Para ello debe instrumentarse el tiempo efectivo de emisión, la cantidad de pulsos, los clics aceptados y el instante usado en el denominador. La prueba debe cubrir el intervalo de \num[round-mode=places,round-precision=2]{\PThreeTransitionLowKm} a \SI[round-mode=places,round-precision=2]{\PThreeTransitionHighKm}{\kilo\meter} y el escenario de visibilidad.

La segunda prioridad es medir una ruta del Campus Miguelete con OTDR o instrumentación equivalente. El resultado debe incluir longitud óptica, atenuación total, conectores, empalmes, fibras disponibles y restricciones de acceso. Con esos datos puede calcularse un presupuesto de enlace específico.

La tercera prioridad es definir una configuración de laboratorio coherente. Deben elegirse modelos exactos de fuente, moduladores, detector y registrador temporal; después se reemplazan los barridos amplios por rangos compatibles con sus hojas de datos y cotizaciones.

La cuarta prioridad es avanzar desde estimadores a protocolo. Esto requiere elegir intensidades dentro de la ejecución, separar estadísticas por base e intensidad, implementar reconciliación, verificación, amplificación de privacidad y autenticación, y aplicar análisis de clave finita.

\section{Cierre}

El trabajo no demuestra que UNSAM pueda desplegar hoy un enlace QKD operativo. Demuestra algo anterior y necesario: cómo construir una simulación auditable, cómo distinguir métricas de desempeño de tasas con interpretación criptográfica y cómo detectar límites antes de comprar hardware. Esa base permite que la siguiente decisión se apoye en datos medidos y no en una extrapolación de curvas.
